\begin{exercisechunk}
\item \textbf{EM for a linear transition with Gaussian noise.}
\label[exercise]{ex:l5-lds-em}
\exerciseprofile{3}{3}{3}{\exproof\quad\excalculation}{%
Derive the exact E- and M-steps for a Gaussian transition from the conditional
moments of the hidden trajectory.}
Use the density form of EM from \cref{ex:l5-continuous-em}. Let
\(S_0,\ldots,S_T\in\R^d\). Relative to Lebesgue measure, suppose the
transition density is
\[
 p_A(s_{t+1}\mid s_t)
 =(2\pi\sigma^2)^{-d/2}
 \exp\!\left(-\frac{\|s_{t+1}-As_t\|^2}{2\sigma^2}\right),
 \qquad t=0,\ldots,T-1,
\]
where \(\sigma^2>0\) is fixed. Let \(p_0\) be the initial-state density
and let \(g_\eta(z_t\mid s_t)\) be the observation density. Assume that
neither depends on \(A\). Denote the joint law by \(Q_{A,\eta}\). Its
complete-data density is
\[
 q_{A,\eta}(s_{0:T},z_{1:T})
 =p_0(s_0)\prod_{t=0}^{T-1}p_A(s_{t+1}\mid s_t)
   \prod_{t=1}^T g_\eta(z_t\mid s_t).
\]
Write its observed-data density as
\[
 m_{A,\eta}(z_{1:T})
 =\int q_{A,\eta}(s_{0:T},z_{1:T})\,ds_{0:T}.
\]
Fix an observed sequence \(z_{1:T,\mathrm{obs}}\), and assume that
\(m_{A_k,\eta_k}(z_{1:T,\mathrm{obs}})>0\) at every iteration. At iteration
\(k\), let
\[
 \Pi_{k,\mathrm{obs}}(ds_{0:T})
 =Q_{A_k,\eta_k}
   (ds_{0:T}\mid Z_{1:T}=z_{1:T,\mathrm{obs}})
\]
be the conditional law of the entire hidden trajectory given the complete
observed sequence. This is often called the \emph{smoothing distribution}.
In contrast, a filtering distribution at time \(t\) conditions only on
observations through time \(t\).

For this fixed iteration and observed sequence, abbreviate the conditional
means, covariances, and adjacent-state cross-covariances by
\begin{align*}
 \mu_t&=\E_{\Pi_{k,\mathrm{obs}}}[S_t],
 &P_t&=\E_{\Pi_{k,\mathrm{obs}}}
   [(S_t-\mu_t)(S_t-\mu_t)^\top],\\
 P_{t,t+1}
 &=\E_{\Pi_{k,\mathrm{obs}}}
   [(S_t-\mu_t)(S_{t+1}-\mu_{t+1})^\top].
\end{align*}
\begin{enumerate}[(a),beginpenalty=10000]
\item Derive the density of \(\Pi_{k,\mathrm{obs}}\) with respect to
Lebesgue measure on the hidden trajectory:
\[
 \pi_{k,\mathrm{obs}}(s_{0:T})
 =\frac{q_{A_k,\eta_k}
   (s_{0:T},z_{1:T,\mathrm{obs}})}
 {m_{A_k,\eta_k}(z_{1:T,\mathrm{obs}})}.
\]
Identify the numerator and normalizing constant.

\item Write one exact EM iteration for \((A,\eta)\): give its E-step and its
M-step as an optimization of the expected complete-data log density under
\(\Pi_{k,\mathrm{obs}}\). Then show that the part of this M-step objective
that depends on \(A\) is
\[
 -\frac{1}{2\sigma^2}\sum_{t=0}^{T-1}
   \E_{\Pi_{k,\mathrm{obs}}}[\|S_{t+1}-AS_t\|^2].
\]
\item Derive the following expression for each summand:
\[
 \E_{\Pi_{k,\mathrm{obs}}}[\|S_{t+1}-AS_t\|^2]
 =\|\mu_{t+1}-A\mu_t\|^2+\operatorname{tr}(P_{t+1})
  -2\operatorname{tr}(A P_{t,t+1})
  +\operatorname{tr}(A^\top A P_t).
\]
\item Assume that
\[
 M_{k,\mathrm{obs}}
 =\sum_{t=0}^{T-1}\E_{\Pi_{k,\mathrm{obs}}}[S_tS_t^\top]
\]
is invertible. Derive the exact M-step update for \(A\), expressed in
terms of \(\mu_t,P_t\), and \(P_{t,t+1}\).
\end{enumerate}
\end{exercisechunk}
